\begin{pastexam}[label={px:nsysu-107-osds-02}]{中山資工所甲組 107 學年度 第 6 題}
（問答、證明）

\textbf{[Advanced Data Structures: 15\%]}
\begin{enumerate}[label=(\arabic*)]
\item Given an $n$-key B-tree with minimum degree $t$, what is the upper bound of tree height?
Prove the correctness of your answer. (8\%)
\item What is a B*-tree? (2\%)
\item Given a binomial tree $B_k$, show that there are exactly $C(k, i)$ nodes at depth $i$, where
$C(k, i)$ denotes a combination function. (5\%)
\end{enumerate}
\end{pastexam}

\begin{pastsoln}
(1)(2) 是 B 樹（資結書第~11.2~節），(3) 是二項樹（資結書第~8.5~節）。

\textbf{(1)} 題目沒有定義高度。 用最小分支度 $t$ 描述 B 樹的是 CLRS，它的高度數邊（只有根的樹高度是
$0$），下面照這個定義作答，最後換成資結書的定義。

答案：$n \ge 1$、$t \ge 2$ 時，高度 $h \le \log_t \dfrac{n + 1}{2}$。（資結書數節點的高度是 $h + 1$，
寫成 $h_{\text{資結書}} \le 1 + \log_t \frac{n+1}{2}$，就是資結書定理~11.1 在 $m = 2t$、$c = t$ 的情況。）

\textbf{要證的是}：任何有 $n \ge 1$ 個鍵、高度 $h$ 的 B 樹都有 $n \ge 2t^h - 1$。 兩邊加 $1$ 除以 $2$
再取 $\log_t$，就得到 $h \le \log_t \frac{n+1}{2}$。 方法是數每一層最少有幾個節點。

用到的定義（CLRS）：根至少 $1$ 個鍵；根以外的節點至少 $t - 1$ 個鍵；有 $k$ 個鍵的內部節點有 $k + 1$
個子節點，所以根如果不是葉節點至少 $2$ 個子節點，根以外的內部節點至少 $t$ 個子節點；所有葉節點的
深度都是 $h$。

令 $N_d$ 是深度 $d$ 的節點數。 \textbf{斷言}：$h \ge 1$ 時，對 $1 \le d \le h$，$N_d \ge 2t^{d-1}$。
對 $d$ 歸納。 $d = 1$：根不是葉節點（$h \ge 1$），至少 $2$ 個子節點，$N_1 \ge 2 = 2t^0$。 設
$1 \le d < h$ 時 $N_d \ge 2t^{d-1}$。 深度 $d$ 的節點不是根，也不是葉節點（葉節點都在深度 $h > d$），
所以各至少有 $t$ 個子節點，而且不同節點的子節點不同，$N_{d+1} \ge t N_d \ge 2t^d$。 斷言成立。

數鍵。 根至少 $1$ 個鍵，深度 $1$ 到 $h$ 的每個節點至少 $t - 1$ 個鍵：
\[
n \;\ge\; 1 + (t - 1) \sum_{d=1}^{h} 2t^{d-1}
\;=\; 1 + 2(t - 1) \cdot \frac{t^h - 1}{t - 1}
\;=\; 2t^h - 1.
\]
$h = 0$ 時只有根，$n \ge 1 = 2t^0 - 1$，也成立。 所以 $t^h \le \frac{n+1}{2}$，
$h \le \log_t \frac{n+1}{2}$。 證畢。

這個上界可以達到：根剛好 $1$ 個鍵、其他每個節點剛好 $t - 1$ 個鍵的 B 樹，上面每個不等號都是等號。
例如 $t = 2$、$h = 2$：根 $1$ 個鍵、$2$ 個子節點、$4$ 片葉子，每個節點 $1$ 個鍵，共 $7 = 2 \cdot 2^2 - 1$
個鍵。

\textbf{(2)} B$^*$ 樹（B$^*$-tree，Knuth 提出）是 B 樹的變形，目的是讓節點更滿、樹更矮：
\begin{itemize}
\item 根以外的每個節點至少要\textbf{三分之二滿}：$m$ 階時至少 $\lceil (2m-1)/3 \rceil$ 個子節點
（B 樹只要求半滿，$\lceil m/2 \rceil$ 個）。
\item 插入時節點滿了，\textbf{先不分裂}，而是把多出來的鍵經由父節點轉給一個還沒滿的相鄰兄弟
（重新分配）。
\item 兄弟也滿了，才把這\textbf{兩個}滿的節點連同父節點裡的分隔鍵，分成\textbf{三個}各約三分之二滿
的節點（2 分 3），父節點多一個鍵。
\end{itemize}
效果是空間使用率從 B 樹的至少約 $50\%$ 提高到至少約 $67\%$，分裂次數變少，同樣的鍵數下樹高不超過
B 樹。 代價是插入要多讀寫兄弟節點，程式比較複雜。 資結書只講 B 樹和 B$^{+}$ 樹，B$^*$ 樹不在書裡；
另外有些書把「B$^{+}$ 樹加上 2 分 3 的分裂規則」也叫 B$^*$ 樹，作答時寫出上面的核心規則即可。

\textbf{(3)} 二項樹的定義（資結書第~8.5~節）：$B_0$ 是一個節點；$B_k$ 由兩棵 $B_{k-1}$ 組成，其中一棵
的根成為另一棵的根最左邊的子節點。 令 $D(k, i)$ 是 $B_k$ 裡深度 $i$ 的節點數（根的深度是 $0$），要證
$D(k, i) = C(k, i)$，$0 \le i \le k$（$i > k$ 或 $i < 0$ 時兩邊都是 $0$）。

對 $k$ 歸納。

\textbf{基底} $k = 0$：$B_0$ 只有根，$D(0, 0) = 1 = C(0, 0)$，其他深度是 $0$。

\textbf{歸納}：設 $D(k-1, i) = C(k-1, i)$ 對所有 $i$ 成立。 $B_k$ 由兩棵 $B_{k-1}$ 組成，叫它們
$T_1$（根就是 $B_k$ 的根）和 $T_2$（根掛在 $T_1$ 的根底下）。
\begin{itemize}
\item $T_1$ 裡的節點，在 $B_k$ 的深度和在 $T_1$ 裡一樣，深度 $i$ 的有 $D(k-1, i)$ 個。
\item $T_2$ 的根變成深度 $1$，$T_2$ 裡每個節點的深度都加 $1$；所以 $B_k$ 裡深度 $i$ 的、來自 $T_2$
的節點，是 $T_2$ 裡深度 $i - 1$ 的節點，有 $D(k-1, i-1)$ 個。
\end{itemize}
兩部分沒有重疊，所以
\[
D(k, i) = D(k-1, i) + D(k-1, i-1) = C(k-1, i) + C(k-1, i-1) = C(k, i),
\]
最後一步是巴斯卡公式。 證畢。

例如 $B_3$ 各深度的節點數是 $1, 3, 3, 1$，正是 $C(3, 0), C(3, 1), C(3, 2), C(3, 3)$。 這也是「二項」
樹名稱的由來。
\end{pastsoln}
